\documentclass[10pt]{article}
\usepackage[T1]{fontenc}
\usepackage[utf8]{inputenc}
\usepackage[french]{babel}
\usepackage{amsmath,amssymb,mathrsfs,mathpazo}
\usepackage{geometry}
\geometry{a4paper,margin=2cm}
\usepackage{graphicx,array,tabularx,tabularray,colortbl,multirow,multicol,booktabs,enumitem,pifont,ulem,fancybox,diagbox,cancel}
\usepackage{pstricks,pstricks-add,pst-plot,pst-tree,pst-node,pst-text,pst-func,pst-math,pst-bspline,pst-3dplot,pst-eucl,pst-solides3d}
\let\pscalculate\undefined
\usepackage{pst-bezier}
\usepackage{tikz,pgfplots,tkz-tab}
\usetikzlibrary{arrows,arrows.meta,patterns,calc,positioning,trees,decorations.markings,shapes.geometric,shapes.misc,shapes.arrows,intersections,decorations.pathmorphing,decorations.pathreplacing,backgrounds,shadows,babel}
\pgfplotsset{compat=1.18}
\usepgfplotslibrary{fillbetween}
\usepackage[np]{numprint}
\usepackage{listings,makecell,pas-tableur,twoopt}
\let\barre\undefined
\usepackage[pstricks]{tablvar}
\definecolor{darkgreen}{RGB}{0,100,0}
\providecommand{\red}{\color{red}}
\providecommand{\blue}{\color{blue}}
\providecommand{\green}{\color{green}}
\providecommand{\black}{\color{black}}
\providecommand{\white}{\color{white}}
\providecommand{\gray}{\color{gray}}
\providecommand{\cyan}{\color{cyan}}
\providecommand{\magenta}{\color{magenta}}
\providecommand{\yellow}{\color{yellow}}
\providecommand{\violet}{\color{violet}}
\providecommand{\orange}{\color{orange}}
\providecommand{\euro}{\texteuro}
\providecommand{\EUR}{\texteuro}
\providecommand{\ieme}{\textsuperscript{e}}
\providecommand{\degres}{\ensuremath{^\circ}}
\providecommand{\degree}{\ensuremath{^\circ}}
\providecommand{\vv}[1]{\overrightarrow{#1}}
\providecommand{\Coord}[1]{#1}
\providecommand{\decosix}{}
\providecommand{\starredbullet}{\ensuremath{\bullet}}
\pagestyle{empty}
\setlength{\parindent}{0pt}
\usepackage{adjustbox}\DeclareUnicodeCharacter{2236}{:}\DeclareUnicodeCharacter{2190}{\ensuremath{\leftarrow}}
\providecommand{\R}{\mathbb{R}}
\providecommand{\N}{\mathbb{N}}
\providecommand{\Z}{\mathbb{Z}}
\providecommand{\Q}{\mathbb{Q}}
\providecommand{\D}{\mathbb{D}}
\providecommand{\C}{\mathbb{C}}
\providecommand{\e}{{\mathrm{e}}}
\providecommand{\vect}[1]{\overrightarrow{#1}}
\providecommand{\vectt}[1]{\overrightarrow{\mathrm{#1}}}
\providecommand{\vvt}[1]{\overrightarrow{\mathrm{#1}}}
\providecommand{\np}[2][]{{#2~#1}}
\providecommand{\vv}[1]{{\overrightarrow{#1}}}
\renewcommand{\d}{{\mathrm{d}}}
\providecommand{\degres}{{^\circ}}
\providecommand{\degre}{{^\circ}}
\providecommand{\blue}{\color{blue}}
\providecommand{\red}{\color{red}}
\providecommand{\enskip}{\quad}
\providecommand{\llbracket}{\lbrack\!\lbrack}
\providecommand{\rrbracket}{\rbrack\!\rbrack}
\providecommand{\euro}{{\text{€}}}
\providecommand{\no}{{\text{n°}}}
\providecommand{\interval}[2]{{[#1;#2]}}
\providecommand{\intervFO}[2]{{[#1;#2[}}
\providecommand{\dotfill}{\ldots}

\providecommand{\ds}{\displaystyle}
\providecommand{\textding}[1]{#1}
\providecommand{\Oij}{\(\mathrm{O};\vec{\imath},\vec{\jmath}\)}
\providecommand{\Oijk}{\(\mathrm{O};\vec{\imath},\vec{\jmath},\vec{k}\)}
\providecommand{\Ouv}{\(\mathrm{O};\vec{u},\vec{v}\)}
\providecommand{\intervFF}[2]{\ensuremath{[#1;#2]}}
\renewcommand{\euro}{\ifmmode\text{\texteuro}\else\texteuro\fi}
\DeclareUnicodeCharacter{00B0}{\ifmmode^{\circ}\else\textdegree\fi}
\DeclareUnicodeCharacter{20AC}{\ifmmode\text{\texteuro}\else\texteuro\fi}
% Énoncés : Éducation nationale. Transcription et corrigés : APMEP, Denis Vergès et contributeurs. Découpage et classement : CoopMaths, Sébastien Lozano, Éric Elter et contributeurs. Adaptation de la présentation : Maths Explorer.
% CC BY-SA : https://coopmaths.fr/www/about/
% https://coopmaths.fr/alea/static/bac/2026/tex/bac_2026_06_sujet1_polynesie_2_cor.tex
\begin{document}
\noindent\textit{APMEP / CoopMaths. CC BY-SA. Présentation adaptée par Maths Explorer.}\par\medskip
\begin{enumerate}
\item Soit $f$ la fonction définie sur $\R$ par $f(x) = (x + 1)\e^{-x}$.

\textbf{Affirmation 1~:} La limite de la fonction $f$ en $+\infty$ est $+\infty$.

\medskip

\begin{tabular}{@{\hspace{0.05\linewidth}} | p{0.92\linewidth}}
D'après les croissances comparées, on sait que $\ds\lim_{x\to +\infty} \dfrac{\e^{x}}{x}=+\infty$ donc $\ds\lim_{x\to +\infty} \dfrac{x}{\e^{x}}=0$, autrement dit $\ds\lim_{x\to +\infty} x\,\e^{-x}=0$.

De plus on sait que $\ds\lim_{x\to +\infty} \e^{-x}=0$.

Par somme; $\ds\lim_{x\to +\infty} x\,\e^{-x} + \e^{-x}=0$
et donc $\ds\lim_{x\to +\infty} f(x)=0$.

\hfill \textbf{Affirmation 1 fausse} 
\end{tabular}

\medskip

\item On considère la fonction $g$ définie sur $\R$ par $g(x) = x\e^x$.

\textbf{Affirmation 2~:} La fonction $g$ est convexe sur l'intervalle $[-2~;~+\infty[$.

\medskip

\begin{tabular}{@{\hspace{0.05\linewidth}} | p{0.92\linewidth}}
Pour déterminer la convexité de la fonction $g$, on étudie le signe de sa dérivée seconde $g''$.

$g(x)=x\e^{x}$ donc $g'(x)= 1\times \e^{x} + x \times \e^{x}=\left (x+1\right )\e^{x}$

$g''(x) = 1 \times \e^{x} + \left (x+1\right )\times \e^{x} = \left (x+2\right )\e^{x}$

Pour tout réel $x$, $\e^{x}>0$, et sur $]-2\;;\; +\infty[$, $x+2>0$.

Donc $g''>0$ sur $]-2\;;\; +\infty[$, donc la fonction $g$ est convexe sur cet intervalle.

\hfill \textbf{Affirmation 2 vraie} 
\end{tabular}

\medskip

\item \textbf{Affirmation 3~:} Une intégration par parties permet d'obtenir~:
$\displaystyle\int_{0}^{\pi} x\sin x\,\text{d}x = -\pi$.

\medskip

\begin{tabular}{@{\hspace{0.05\linewidth}} | p{0.92\linewidth}}
Sur l'intervalle $]0\;;\;\pi[$, $\sin\,x>0$ donc $x\,\sin x >0$;

D'après les propriétés de l'intégrale, on en déduit que 
$\displaystyle\int_{0}^{\pi} x\sin x\,\text{d}x>0$.

\hfill \textbf{Affirmation 3 fausse} 
\end{tabular}

\medskip

\item On considère l'équation différentielle $(E)$ : $y' = 2y - 5$ et la fonction $h$ définie sur $\R$ par $h(x) = -\dfrac{3}{2}\e^{2x} + \dfrac{5}{2}$.

\textbf{Affirmation 4~:} La fonction $h$ est la solution de $(E)$ qui vaut $1$ en $0$.

\medskip

\begin{tabular}{@{\hspace{0.05\linewidth}} | p{0.92\linewidth}}
 $h(x) = -\dfrac{3}{2}\e^{2x} + \dfrac{5}{2}$
 donc $h'(x)= -\dfrac{3}{2}\times 2\e^{2x}=-3\e^{2x}$
 
 \begin{list}{\textbullet}{}
 \item $2h(x)-5 = 2\left (-\dfrac{3}{2}\e^{2x}+\dfrac{5}{2}\right )-5 = -3\e^{2x}+5-5=-3\e^{2x}=h'(x)$
 
 Donc la fonction $h$ est solution de l'équation différentielle $(E)$.
 
 \item $h(0)=-\dfrac{3}{2}\e^{0}+\dfrac{5}{2}=-\dfrac{3}{2}+ \dfrac{5}{2}=\dfrac{2}{2}=1$
 \end{list}

\hfill \textbf{Affirmation 4 vraie} 
\end{tabular}

\medskip

\item On considère la fonction \texttt{mystere} définie ci-dessous en Python~:

\begin{center}












\begin{tabular}{|l|}\hline
\texttt{def mystere(n):}\\
\quad \texttt{ f = 1}\\
\quad \texttt{ for k in range(1, n+1):}\\
\quad \quad \texttt{f = f * k}\\
\quad  \texttt{return f}\\ \hline
\end{tabular}

\end{center}

\textbf{Affirmation 5~:} Le nombre renvoyé lors de l'exécution de \texttt{mystere(5)} est $120$.


\medskip

\begin{tabular}{@{\hspace{0.05\linewidth}} | p{0.92\linewidth}}
On fait tourner cet algorithme: comme on lance \texttt{mystere(5)}, la variable \texttt{n} vaut 5 et on fait tourner la boucle \texttt{for} pour \texttt{k} variant entre 1 et 5.

\begin{center}
\begin{tabular}{| *6{c |} }
\hline
\texttt{n} & 1&2&3&4&5\\
\hline
\texttt{f} & $1\times 1=1$ & $1\times 2=2$ & $2\times 3=6$ & $6 \times 4=24$ & $24\times 5=120$\\
\hline
\end{tabular}
\end{center}

Pour \texttt{n=5} on obtient \texttt{f=120}; la fonction \texttt{mystere} calcule la factorielle.

\hfill \textbf{Affirmation 5 vraie} 
\end{tabular}

\end{enumerate}

\end{document}
